\chapter{Discussion}\label{ch:discussion}
This chapter presents the key takeaways of this research and outlines directions for future work. I first discuss how orthogonal statechart modelling for SoTS enables controlled and adaptive behaviour under disturbance (\secref{sec:discussion-utility}), then examine its implications for SoTS engineering practice (\secref{sec:discussion-implications}), and finally outline directions for prospective researchers (\secref{sec:discussion-future}).

\section{On the Utility of Model-Driven Coordination in SoTS}\label{sec:discussion-utility}
The results in \secref{sec:evaluation} show that reliability in SoTS can be achieved by controlling how constituents participate in the system, rather than relying solely on service-level fault tolerance mechanisms.

\subsection{Model-Driven Coordination and Control}
The results in \secref{sec:evaluation} demonstrate that coordination behaviour can be defined and controlled through lifecycle statechart models, rather than implemented within individual system components. This operationalizes the state-based reliability model introduced in \chref{ch:reliable-sots}, where disturbances arise from interactions between belonging and health.  In the demonstrated architecture (\secref{sec:architecture}), executable statecharts act as a coordination layer that controls when and how constituents contribute to system-level behaviour. Participation is determined by constraints encoded in the lifecycle model, ensuring that only constituents in valid states contribute to the event stream and CEP processing. This prevents disturbance-causing configurations by restricting how degraded or unreliable constituents participate. This externalizes coordination to the system level, enabling behaviour to be specified in models and controlled uniformly at runtime. As a result, it reduces implementation complexity and allows reliability to be expressed and maintained as constraints on system behaviour.

\subsection{Decoupling Coordination from Service Logic}

The proposed approach separates constituent internals from system-level concerns by externalising SoTS coordination into lifecycle models. DT constituents participate through defined states and transitions without exposing internal behaviour or implementation details. Behaviour is controlled by a shared model rather than embedded within individual components, allowing constituents to evolve independently while maintaining system-level consistency. This aligns with operational and managerial independence in SoS \cite{maier1998architecting} and architectural principles that limit exposure to stakeholder-relevant concerns \cite{iso42010}.

This separation also reduces information exposure and integration complexity. Interaction is mediated through lifecycle states, so internal models do not need to account for external system dynamics. This aligns with privacy perspectives such as the NIST Privacy Framework \cite{enterprise2020nist}, where risks are addressed across the data lifecycle. By constraining elligibility in the SoTS based on lifecycle state, the model limits unnecessary data flows and unintended exposure, which is critical in SoS environments where independently managed systems interact \cite{olivero2024systematic}. It also clarifies engineering responsibilities: SoTS-level developers define system-level constraints and guarantees, while service-level developers focus on local behaviour.


\subsection{Reducing Complexity with Model-to-Code Generation}

% \begin{table}[t]
% \centering
% \footnotesize
% \caption{Code complexity across reliability levels}
% \begin{tabular}{
% >{\arraybackslash}p{3cm}
% >{\centering\arraybackslash}p{2cm}
% >{\centering\arraybackslash}p{2.6cm}
% >{\centering\arraybackslash}p{2.6cm}
% >{\centering\arraybackslash}p{2.6cm}
% }
% \toprule
% \textbf{Reliability Level} 
% & \textbf{Lines of Code} 
% & \textbf{Avg Cyclomatic Complexity} 
% & \textbf{Max Cyclomatic Complexity} 
% & \textbf{Programming Time (s)} \\
% \midrule
% Level 1 Reactive   & 942  & 1.86 & 14 & 663 \\
% Level 2 Proactive  & 1169 & 2.07 & 18 & 1143 \\
% Level 3 Controlled & 1343 & 2.00 & 21 & 1512 \\
% Level 4 Adaptive   & 1469 & 2.15 & 23 & 1462 \\
% \bottomrule
% \end{tabular}
% \end{table}


\begin{table}[htp]
\centering
\footnotesize
\caption{Code complexity across reliability levels}
\label{tab:complexity}
\begin{tabular}{
>{\raggedright\arraybackslash}p{3.25cm}
>{\centering\arraybackslash}p{2cm}
>{\centering\arraybackslash}p{3cm}
>{\centering\arraybackslash}p{3cm}
}
\toprule
\textbf{SoTS Reliability Level} 
& \textbf{Lines of Code} 
& \textbf{Cyclomatic Complexity} 
& \textbf{Programming Time (min)} \\
\midrule
Level 1 Reactive   & 942  & 14 & 11 \\
Level 2 Proactive  & 1169 & 18 & 19 \\
Level 3 Controlled & 1343 & 21 & 25 \\
Level 4 Adaptive   & 1469 & 23 & 24.5 \\
\bottomrule
\end{tabular}
\end{table}

\id{Start by tying this to the main idea, i.e., first sentence should be that model-driven coordination enables this + explain how. Otherwise this doesn't belong in this subsection.}

In the demonstration (\secref{sec:itemis-create}), Itemis CREATE is used to generate executable code directly from lifecycle statechart models. \tabref{tab:complexity} summarizes code complexity metrics for the generated implementations across each SoTS reliability level (\secref{sec:reliability-levels}). These metrics were obtained using the Radon analysis library\footnote{\url{https://radon.readthedocs.io/en/latest/intro.html}}. They include lines of code (LOC) and cyclomatic complexity (number of independent execution paths, indicating control flow complexity \cite{coleman1994using, oman1992metrics}). From Level~1 to Level~4, LOC increases from 942 to 1469 (56\%). Cyclomatic complexity increases from 14 to 23 (64\%).
% study allows us to show the tendency of things

Normally, the logic required to handle these behaviours would be distributed across components. This requires manual implementation and increases the risk of inconsistencies. A model-driven approach instead defines constraints directly in statecharts. It implements them through generated code. This makes the logic easier to reason about, modify, and verify. It also reduces duplication and development effort.




\section{Future Directions}\label{sec:discussion-future}
This section identifies opportunities to extend the proposed approach to support additional concerns and more complex system behaviour. Lifecycle-based coordination provides a flexible foundation that can be expanded to incorporate additional system-level constraints, adaptive strategies, and formal analysis techniques.

\subsection{Modelling Additional Concerns as Orthogonal Regions}

\secref{sec:two-dimensional-disturbances} introduced a two-dimensional model of disturbances based on belonging and health. While these dimensions support reasoning about reliability, they do not capture all factors that influence participation. Additional concerns such as trust, authorization level, confidentiality, and data integrity also affect whether and how DTs interact \cite{li2023trust, kuruppuarachchi2023trusted}. These concerns are well recognized in distributed DT settings, where issues such as data sharing and access control directly affect which DTs are suitable for collaboration \cite{tripathi2024stakeholders}. They can be incorporated into the model as orthogonal regions, where each region captures a distinct aspect influencing participation. This allows eligibility to be determined based on multiple system-level factors rather than belonging and health alone. However, expanding the model introduces additional complexity. Stakeholders must agree on how these concerns are defined and how they influence participation \cite{carney2005topics}. While this enables more precise control over contributions, it requires careful reasoning about interactions between regions and extends the model beyond reliability to broader dependability concerns \cite{avizienis2004basic}.

\subsection{Dynamic Fault Tolerance in SoTS}
At the highest reliability level (\secref{sec:lv4}), the model allows constituents to continue participating under a restricted role when their condition degrades. This is realized in \secref{sec:demo-compensation} through a compensation component that estimates missing data. Different fault tolerance strategies can be applied in these situations \cite{hanmer2013patterns}, each introducing trade-offs in resource usage and predictability \cite{narasimhan2002trade-offs}. In dynamic environments, selecting strategies at runtime allows the system to respond more effectively than fixed approaches \cite{desouza2022systematic, ballesteros2022infrastructure}.

The proposed model supports this by introducing an additional orthogonal region representing the selection of fault tolerance strategies. States in this region represent different strategies, while transitions are influenced by changes in other regions, enabling adaptive behaviour. For example, estimation can handle intermittent missing data, while more severe degradation can trigger data fusion using observations from nearby DTs. Because statecharts are executable, learning-based methods such as reinforcement learning can be used to refine transitions and guards based on observed system performance. This enables the model to improve how fault tolerance strategies are selected \cite{zhang2020fault, dong2023deep}.


\subsection{Formal Verification of SoTS}
Through adopting a state based modelling approach to controlling SoTS (\secref{sec:statecharts-reliability}), we gain the ability to formally verify system behaviour. Models allow system behaviour to be represented explicitly, enabling formal verification through systematic analysis of all possible system executions. It becomes possible to reason about how the system can evolve under all different conditions. This allows properties such as safety (forbidden configurations are never reached), liveness (desired progression occurs), and invariants over belonging and health to be checked with guarantees \cite{henzinger2023quantitative}. This can be supported through the integration of formal analysis tools into the modelling workflow, enabling automated verification as the system evolves \cite{drusinsky2011modeling}.

Further analysis can be supported through model-driven engineering (MDE) techniques. Behaviour is defined once in the statechart and systematically transformed into other models for analysis \cite{biehl2010literature}. These transformations provide alternative views of the same behaviour, making properties visible that are not explicit in the original model \cite{bezivin2006model}. For example, while statecharts capture the lifecycle of constituents, they do not directly expose resource usage, synchronization structure, or global system conditions. By transforming the model into representations such as Petri nets \cite{cortadella1995synthesizing, hei2011automatic}, these aspects can be analysed explicitly, enabling detection of deadlocks, unreachable states, and resource contention. Insights from this analysis can be fed back into the statechart through additional constraints or guarded transitions. This strengthens the link between execution and verification, supporting more reliable SoTS behaviour.
